% titel: Analogie Translation und Rotation
% skript: Kap. Rotation, Analogie Translation-Rotation

Zwischen geradliniger Bewegung (Translation) und Drehbewegung (Rotation) um
eine feste Achse besteht eine vollständige Analogie: Jeder Translationsgrösse
entspricht eine Rotationsgrösse, jedem Translationsgesetz ein
Rotationsgesetz. Die Tabelle stellt die einander entsprechenden Grössen
gegenüber.

\begin{werkzeugkasten}[Analogie-Tabelle]
\begin{center}
\renewcommand{\arraystretch}{1.35}
\begin{tabular}{lll}
\toprule
\textbf{Translation} & \textbf{Grösse} & \textbf{Rotation}\\
\midrule
Ort / Winkel & $x$ & $\varphi$ \\
Geschwindigkeit & $v=\dot x$ & $\omega=\dot\varphi$ \\
Beschleunigung & $a=\dot v$ & $\alpha=\dot\omega$ \\
Masse / Trägheitsmoment & $m$ & $I$ \\
Kraft / Drehmoment & $F$ & $M=\tau$ \\
Impuls / Drehimpuls & $p=mv$ & $L=I\omega$ \\
kinetische Energie & $\frac12 mv^2$ & $\frac12 I\omega^2$ \\
\bottomrule
\end{tabular}
\end{center}

\textbf{Gesetze in Analogie:} $F=ma$ entspricht $M=I\alpha$; der Impulssatz
$F=\dot p$ dem Drehimpulssatz $M=\dot L$; die Energie
$\int F\,dx$ dem $\int M\,d\varphi$.
\end{werkzeugkasten}

\begin{werkzeugkasten}[Nutzen]
Mit der Analogie kann man ein bekanntes Translationsproblem direkt auf die
Rotation übertragen: Bewegungsgleichungen, Erhaltungssätze, Energie- und
Drehimpulsbetrachtungen haben dieselbe Struktur.
\end{werkzeugkasten}

\begin{achtung}[Hinweis]
\begin{itemize}
  \item Bei der Rotation um eine feste Achse sind alle Grössen Skalare mit
       Vorzeichen; beim vollständigen (Vektor-)Problem übernehmen $\vec\tau$
       und $\vec L$.
  \item $I$ hängt von der Massenverteilung und Achse ab (Satz von Steiner),
       $m$ ist invariant.
\end{itemize}
\end{achtung}
